% R1的原HEFT适配版H/R信息方案：取H02逐槽已保存调度，按照同一真实H模型画AB场景。
% blocks的真实时长为stage_catalog原H响应和合法OPP的确定值；不显示候选R预测时长。
\begin{figure}[htbp]
\centering
\resizebox{0.98\textwidth}{!}{%
\begin{tikzpicture}[x=.089cm,y=1cm,font=\small,>=Latex]
 \node[anchor=west,font=\bfseries] at (-13,8.05) {H信息规划：真实H能量0.881088 J};
 \node[anchor=west,font=\bfseries] at (-13,4.15) {R信息规划：真实H能量0.917359 J};
 \foreach \xx in {0,25,50,75,100,125}{
  \draw[densely dotted,gray!50] (\xx,.78)--(\xx,7.65);
  \node[below,font=\scriptsize] at (\xx,.65) {\xx};
 }
 \draw[->,thick] (0,.60)--(132,.60) node[right,font=\scriptsize] {时间/ms};
 \foreach \yy/\lab in {7.25/CPU1,6.52/CPU2,5.79/GPU0,5.06/NPU0,3.36/CPU1,2.63/CPU2,1.90/GPU0,1.17/NPU0}{
   \node[anchor=east,font=\scriptsize] at (-4,\yy) {\lab};
   \draw[gray!30] (0,\yy-.26)--(125,\yy-.26);
 }
 \filldraw[draw=blue!70!black,fill=blue!18,thick] (70,4.87) rectangle (112,5.25);
 \node[font=\tiny] at (91,5.06) {A.main/L};
 \filldraw[draw=blue!70!black,fill=blue!18,thick] (115,7.06) rectangle (121.5,7.44);
 \node[font=\tiny,above] at (118.25,7.40) {A.post};
 \filldraw[draw=blue!70!black,fill=blue!18,thick] (5,6.33) rectangle (19,6.71);
 \node[font=\tiny] at (12,6.52) {A.pre/L};
 \filldraw[draw=orange!75!black,fill=orange!18,thick] (25,4.87) rectangle (67,5.25);
 \node[font=\tiny] at (46,5.06) {B.main/M};
 \filldraw[draw=orange!75!black,fill=orange!18,thick] (75,7.06) rectangle (81.6,7.44);
 \node[font=\tiny,above] at (78.3,7.40) {B.post};
 \filldraw[draw=orange!75!black,fill=orange!18,thick] (0,7.06) rectangle (22.5,7.44);
 \node[font=\tiny] at (11.25,7.25) {B.pre/L};
 \filldraw[draw=blue!70!black,fill=blue!18,thick] (50,0.98) rectangle (92,1.36);
 \node[font=\tiny] at (71,1.17) {A.main/L};
 \filldraw[draw=blue!70!black,fill=blue!18,thick] (110,3.17) rectangle (116.5,3.55);
 \node[font=\tiny,above] at (113.25,3.51) {A.post};
 \filldraw[draw=blue!70!black,fill=blue!18,thick] (20,2.44) rectangle (34,2.82);
 \node[font=\tiny] at (27,2.63) {A.pre/L};
 \filldraw[draw=orange!75!black,fill=orange!18,thick] (15,0.98) rectangle (50,1.36);
 \node[font=\tiny] at (32.5,1.17) {B.main/H};
 \filldraw[draw=orange!75!black,fill=orange!18,thick] (55,3.17) rectangle (61.6,3.55);
 \node[font=\tiny,above] at (58.3,3.51) {B.post};
 \filldraw[draw=orange!75!black,fill=orange!18,thick] (0,3.17) rectangle (15,3.55);
 \node[font=\tiny] at (7.5,3.36) {B.pre/H};
 \fill[blue!18,draw=blue!65!black] (80,-.05) rectangle (86,.17);
 \node[anchor=west,font=\scriptsize] at (87,.06) {蓝色：A};
 \fill[orange!18,draw=orange!70!black] (110,-.05) rectangle (116,.17);
 \node[anchor=west,font=\scriptsize] at (117,.06) {橙色：B};
\end{tikzpicture}}
\caption{H02中R1的A+B场景两种规划安排（时间块按真实H响应绘制，系统能量同时计入A+C场景）}
\label{fig:ch3-h02-gantt}
\end{figure}
